\documentclass[10pt,a4paper]{amsart}
\usepackage[margin=19mm]{geometry}
\usepackage{amsmath,amssymb,mathtools}
\usepackage[scr=rsfs,cal=euler]{mathalfa}
\usepackage{xfrac}
\usepackage[unicode,hidelinks,bookmarks=false]{hyperref}
\newcommand{\ie}{i.e.\ }
\newcommand{\cf}{cf.\ }
\newcommand{\resp}{respectively\ }
\newcommand{\Subsection}{\subsection}
\providecommand{\nobreakdash}{\nobreak\hbox{-}}
\setlength{\emergencystretch}{2em}
\multlinegap0pt
\usepackage{bm}
\usepackage{bbm}
\usepackage{dsfont}
\usepackage{xspace}
\usepackage{cancel}
\usepackage{mathtools}
\usepackage{amsthm}
\makeatletter
\@ifundefined{thrm}{\newtheorem{thrm}{Thrm}}{}
\@ifundefined{prpstn}{\newtheorem{prpstn}[thrm]{Proposition}}{}
\makeatother
\title[On the Hardy-Ramanujan Theorem]{On the Hardy-Ramanujan Theorem}
\author{Benjamin Durkan}
\date{Source: arXiv:2310.14760v2 (CC BY 4.0)}
\begin{document}
\maketitle

\noindent\emph{Source and licence.} On the Hardy-Ramanujan Theorem. Version arXiv:2310.14760v2 (CC BY 4.0), distributed under the Creative Commons Attribution 4.0 International licence, as recorded in the arXiv licence metadata for this exact version and shown on the abstract page https://arxiv.org/abs/2310.14760v2. Full rights notice: licenses/arxiv-2310.14760-CC-BY-4.0.txt.\par\smallskip

\noindent\textit{Editorial scope.} Counted unit: the Prpstn whose source label is \texttt{prop:gaussian\_Omega} in the article (source0.tex lines 481--486, source label \texttt{prop:gaussian\_Omega}), delivered together with the article's own complete original proof of it (source0.tex lines 487--525, one \texttt{proof} environment of 39 lines). No mathematics is written, repaired or re-derived: statement and proof are the article's own text line for line. Required context retained uncounted: eqn:omega\_upper\_window (lines 399--402); prop:R\_tail (lines 449--454). Every definition, display and label of those lines is retained unchanged. Omitted from this package and not redistributed: 1--398, 403--448, 455--480, 526--933. Nothing inside a retained excerpt is omitted or altered: every retained body is delivered exactly as the article's lines are written, so no omission receipt of any kind accompanies this package. Citations: the retained excerpts carry no citation command at all, so no bibliography entry of the article is transcribed and no reference is redistributed. Reference adaptation: none was needed and none was made; every reference inside the delivered unit resolves inside it, so no unresolved reference or citation is produced. Printed slips kept literal: no printed word, symbol, hypothesis or number of the counted statement or of its proof was altered. Layout and class: the article's own class is \texttt{amsart} with packages including amsmath, amsthm, amsfonts, amssymb, txfonts, graphicx, bm, url, verbatim, siunitx, hyperref; the wrapper is the lane's pinned \texttt{amsart} editorial wrapper at 10pt on A4, which restates the theorem-like environments the retained text needs and adds the article's own macro definitions that the retained text needs (none), with extra wrapper packages (bm, bbm, dsfont, xspace, cancel, mathtools). The delivered numbering is the unit's own. Assets: no retained span contains a figure environment, a graphics-inclusion command, a PDF-page inclusion, a table or a TikZ picture; no image file is copied into this package, the package declares no asset dependency and no image file is redistributed with it. Licence: version arXiv:2310.14760v2 of this article is distributed under the Creative Commons Attribution 4.0 International licence, as recorded in the arXiv licence metadata for this exact version and shown on the abstract page https://arxiv.org/abs/2310.14760v2. A full rights notice is delivered at licenses/arxiv-2310.14760-CC-BY-4.0.txt. Characters: every retained excerpt is pure ASCII, so the delivered engine, XeLaTeX with its default Latin Modern text font, reports no missing character.
    \begin{align}
        \frac{1}{x}\#\{2\le n\le x:\omega(n)-L>A\sqrt{L}\}&\le\frac{8Q}{5A}\exp\left(-\frac{A^2}{2}+\frac{5.096A}{\sqrt{L}}+\frac{2.048A^2}{L}+\frac{A^3}{6\sqrt{L}}\right),\label{eqn:omega_upper_window}\\
        \frac{1}{x}\#\{2\le n\le x:L-\omega(n)>A\sqrt{L}\}&\le\frac{Q}{A}e^{-A^2/2}.\label{eqn:omega_lower_window}
    \end{align}

\begin{prpstn}\label{prop:R_tail}
    For all $x\ge 1$ and all integers $T\ge 0$,
    \begin{equation}\label{eqn:R_tail}
        \frac{1}{x}\#\{n\le x:R(n)\ge T\}\le\frac{11}{2}2^{-T}.
    \end{equation}
\end{prpstn}

\begin{prpstn}\label{prop:gaussian_Omega}
    If $L=\log\log x$, $L\ge 200$ and $1\le A\le\sqrt{L}/12$, then
    \begin{equation}\label{eqn:Omega_window}
        \frac{1}{x}\#\{2\le n\le x:|\Omega(n)-L|>A\sqrt{L}\}\le\frac{332}{A}\exp\left(-\frac{A^2}{2}+\frac{2.29A^3}{\sqrt{L}}\right).
    \end{equation}
\end{prpstn}

\begin{proof}
    For $1\le A\le 3$, the right side of \eqref{eqn:Omega_window} is $>1$, so the estimate is trivial. Assume $A>3$. Define
    $$T_A=\left\lceil\frac{A^2/2+\log(11eA/2)}{\log 2}\right\rceil.$$
    By Proposition \ref{prop:R_tail},
    \begin{equation}\label{eqn:R_tail_TA}
        \frac{1}{x}\#\{n\le x:R(n)\ge T_A\}\le\frac{1}{eA}e^{-A^2/2}.
    \end{equation}
    We shall use
    \begin{equation}\label{eqn:TA_bound}
        T_A\le\frac{3}{2}A^2\qquad (A\ge 3).
    \end{equation}
    Indeed, after accounting for the ceiling, \eqref{eqn:TA_bound} follows from
    $$\log(11eA/2)+\log 2\le\left(\frac{3}{2}\log 2-\frac{1}{2}\right)A^2,$$
    which holds at $A=3$ and then follows by monotonicity. Since $A\le\sqrt{L}/12$, \eqref{eqn:TA_bound} gives
    $$\frac{T_A}{\sqrt{L}}\le\frac{A}{8}.$$
    Put
    $$B=A-\frac{T_A}{\sqrt{L}}.$$
    Then $B\ge 7A/8>1$ and $B\le A\le\sqrt{L}/12$.

    The lower tail is immediate from $\Omega(n)\ge\omega(n)$:
    $$\frac{1}{x}\#\{2\le n\le x:L-\Omega(n)>A\sqrt{L}\}\le\frac{Q}{A}e^{-A^2/2}.$$
    For the upper tail, if $\Omega(n)-L>A\sqrt{L}$, then either $R(n)\ge T_A$, or
    $$\omega(n)-L>A\sqrt{L}-T_A=B\sqrt{L}.$$
    Using \eqref{eqn:omega_upper_window} with $B$ in place of $A$, together with \eqref{eqn:R_tail_TA}, gives
    $$\frac{1}{x}\#\{2\le n\le x:\Omega(n)-L>A\sqrt{L}\}\le\frac{8Q}{5B}e^{E(B)}+\frac{1}{eA}e^{-A^2/2},$$
    where
    $$E(B)=-\frac{B^2}{2}+\frac{5.096B}{\sqrt{L}}+\frac{2.048B^2}{L}+\frac{B^3}{6\sqrt{L}}.$$
    Since $B\ge 7A/8$, the prefactor is at most $64Q/(35A)$. Also, using $A>3$, $L\ge 200$, $B\le A$ and \eqref{eqn:TA_bound},
    \begin{align*}
        E(B)&\le -\frac{A^2}{2}+\frac{3}{2}\frac{A^3}{\sqrt{L}}+\frac{5.096}{9}\frac{A^3}{\sqrt{L}}+\frac{2.048}{3\sqrt{200}}\frac{A^3}{\sqrt{L}}+\frac{1}{6}\frac{A^3}{\sqrt{L}}\\
        &\le -\frac{A^2}{2}+\frac{2.29A^3}{\sqrt{L}}.
    \end{align*}
    Adding the lower tail, the shifted upper tail and \eqref{eqn:R_tail_TA}, we obtain
    \begin{align*}
        \frac{1}{x}\#\{2\le n\le x:|\Omega(n)-L|>A\sqrt{L}\}&\le\left(\frac{Q}{A}+\frac{64Q}{35A}+\frac{1}{eA}\right)\\
        &\quad\times\exp\left(-\frac{A^2}{2}+\frac{2.29A^3}{\sqrt{L}}\right).
    \end{align*}
    Since $Q<117.20$, the coefficient is $<332/A$.
\end{proof}

\end{document}
